\chapter{Atomes polyélectroniques et symboles de termes}\label{ch:b3:many-electron-atoms}

En 1868, une raie jaune du spectre du Soleil révéla un élément inconnu
sur Terre, l'hélium. Quand on l'isola enfin au laboratoire, en 1895, son
spectre ressemblait à celui de \emph{deux} éléments~: deux familles de
raies, attribuées un temps à l'«~orthohélium~» et au «~parahélium~», et
qui n'échangeaient jamais de lumière entre elles. Il n'existe qu'un seul
hélium. Les deux familles sont les \omterm{def:b3:many-electron-atoms:singlet-triplet}{états triplets} et singulets du même
atome, tenus séparés par le spin de l'électron et par une règle plus
fondamentale que toute force~: la fonction d'onde de deux électrons doit
changer de signe quand on les échange. Ce chapitre suit cette règle du
principe de Pauli jusqu'aux \omterm{def:b3:many-electron-atoms:term}{symboles de termes} par lesquels les
spectroscopistes désignent chaque état de chaque atome et de chaque ion
--- y compris les ions des métaux de transition dont les couleurs sont
expliquées au \cref{ch:b3:complex-spectra-magnetism}.

\begin{recall}
Le volume de première année a désigné les électrons par quatre nombres
quantiques $n$, $l$,
$m_l$, $m_s$, a construit les configurations fondamentales avec le
principe de Pauli, la règle de Klechkowski et la règle de Hund, et a
compté les électrons célibataires. Le volume de deuxième année a donné
les énergies des orbitales et les règles de Slater pour l'écrantage. Le
\Cref{ch:b3:quantum-model-systems} a fourni les \omterm{def:b3:quantum-model-systems:operator}{opérateurs}, les \omterm{def:b3:quantum-model-systems:operator}{valeurs
propres}, le moment cinétique du rotateur ($\hat L^2$, de \omterm{def:b3:quantum-model-systems:operator}{valeurs propres}
$l(l+1)\hbar^2$) et l'atome d'hydrogène.
\end{recall}

\section{Spin et spin-orbitales}

L'électron porte un moment cinétique intrinsèque, son spin, de nombre
quantique $s = \frac12$~: $\hat S^2$ a l'unique \omterm{def:b3:quantum-model-systems:operator}{valeur propre} $s(s+1)\hbar^2 =
\frac34\hbar^2$, et $\hat S_z$ les deux \omterm{def:b3:quantum-model-systems:operator}{valeurs propres} $m_s\hbar = \pm\frac12\hbar$.
Les deux fonctions de spin s'écrivent $\alpha$ ($m_s = +\frac12$) et $\beta$
($m_s = -\frac12$)~; elles sont orthonormées, $\langle\alpha|\alpha\rangle =
\langle\beta|\beta\rangle = 1$ et $\langle\alpha|\beta\rangle = 0$. Le spin
n'a pas d'équivalent classique~; ce n'est pas la rotation d'une petite
sphère, et il n'intervient pas du tout dans un hamiltonien non
relativiste.

\begin{definition}[Spin-orbitale]\label{def:b3:many-electron-atoms:spin-orbital}
Une \emph{spin-orbitale}\index{spin-orbitale} est le produit d'une orbitale
spatiale
$\phi(\mathbf r)$ et d'une fonction de spin~: $\phi\alpha$ ou $\phi\beta$.
Une spin-orbitale contient la description complète d'un électron.
\end{definition}

Une orbitale atomique repérée par $(n, l, m_l)$ dans le volume de première
année donne donc deux \omterm{def:b3:many-electron-atoms:spin-orbital}{spin-orbitales}, et c'est pourquoi une sous-couche
contient $2(2l+1)$ électrons. La vraie question est de savoir pourquoi
une \omterm{def:b3:many-electron-atoms:spin-orbital}{spin-orbitale} en contient au plus un.

\section{Électrons indiscernables et déterminant de Slater}

\begin{definition}[Particules indiscernables, fonction d'onde antisymétrique]\label{def:b3:many-electron-atoms:indistinguishable}
Des particules sont \emph{indiscernables}\index{particules indiscernables} si
aucune mesure ne permet de les distinguer~: tout électron est identique à
tout autre. Une fonction d'onde de plusieurs électrons est une
\emph{fonction d'onde antisymétrique}\index{fonction d'onde antisymetrique@fonction d'onde antisymétrique} si elle change de signe quand on
échange les coordonnées (d'espace et de spin) de deux électrons
quelconques~:
$\Psi(\dots,i,\dots,j,\dots) = -\Psi(\dots,j,\dots,i,\dots)$.
\end{definition}

L'indiscernabilité seule impose que $|\Psi|^2$ soit inchangé par un
échange, si bien que $\Psi$ peut au plus changer de signe. La nature a
choisi~: la fonction d'onde de tout système d'électrons --- de tous
fermions --- est antisymétrique. C'est un postulat de la mécanique
quantique, confirmé par tous les spectres atomiques.

\begin{definition}[Déterminant de Slater]\label{def:b3:many-electron-atoms:slater}
Pour $N$ électrons dans $N$ \omterm{def:b3:many-electron-atoms:spin-orbital}{spin-orbitales} différentes $\chi_1, \dots, \chi_N$, le
\emph{déterminant de Slater}\index{determinant de Slater@déterminant de Slater} est
\[
\Psi = \frac{1}{\sqrt{N!}}\begin{vmatrix}
\chi_1(1) & \chi_2(1) & \cdots & \chi_N(1)\\
\chi_1(2) & \chi_2(2) & \cdots & \chi_N(2)\\
\vdots & & & \vdots\\
\chi_1(N) & \chi_2(N) & \cdots & \chi_N(N)
\end{vmatrix},
\]
où la ligne $i$ contient l'électron $i$ dans chaque \omterm{def:b3:many-electron-atoms:spin-orbital}{spin-orbitale}. On
l'écrit en abrégé
$|\chi_1\chi_2\dots\chi_N|$.
\end{definition}

\begin{theorem}[Le principe de Pauli]\label{thm:b3:many-electron-atoms:pauli}
Un \omterm{def:b3:many-electron-atoms:slater}{déterminant de Slater} est antisymétrique, et il s'annule si deux
électrons occupent la même \omterm{def:b3:many-electron-atoms:spin-orbital}{spin-orbitale}. Par conséquent, deux électrons
d'un atome ne peuvent pas avoir les mêmes quatre nombres quantiques.
\end{theorem}

\begin{proof}
Échanger les électrons $i$ et $j$ échange deux lignes du déterminant, ce
qui change son signe~: l'antisymétrie vaut pour toute paire. Si deux
\omterm{def:b3:many-electron-atoms:spin-orbital}{spin-orbitales} sont égales, deux colonnes sont égales et le déterminant
est nul~: un tel état n'existe pas. Deux électrons du même atome ayant
les mêmes $n$, $l$,
$m_l$ et $m_s$ occuperaient la même \omterm{def:b3:many-electron-atoms:spin-orbital}{spin-orbitale}.
\end{proof}

\begin{example}[L'état fondamental de l'hélium]\label{ex:b3:many-electron-atoms:helium-ground}
Avec les deux électrons dans $1s$,
\[
\Psi = \frac{1}{\sqrt2}\begin{vmatrix}1s\alpha(1) & 1s\beta(1)\\ 1s\alpha(2) &
1s\beta(2)\end{vmatrix} = 1s(1)\,1s(2)\cdot\frac{1}{\sqrt2}\big[\alpha(1)\beta(2)
- \beta(1)\alpha(2)\big].
\]
La partie spatiale est symétrique, la partie de spin antisymétrique~: les
deux spins sont appariés, de spin total nul.
\end{example}

\section{L'hélium~: répulsion coulombienne et échange}

Excitons un électron de l'hélium vers $2s$. On peut construire quatre
déterminants à partir de
$1s\alpha$, $1s\beta$, $2s\alpha$, $2s\beta$, avec un électron dans chaque
orbitale. Réarrangés, ils deviennent des produits d'une fonction spatiale
et d'une fonction de spin~:
\[
\Psi_\pm = \frac{1}{\sqrt2}\big[1s(1)2s(2) \pm 2s(1)1s(2)\big]\times(\text{fonction
de spin}).
\]
La fonction spatiale antisymétrique $\Psi_-$ doit être associée à une
fonction de spin symétrique, et il y en a trois~: $\alpha(1)\alpha(2)$,
$\beta(1)\beta(2)$ et $\frac{1}{\sqrt2}[\alpha(1)\beta(2) + \beta(1)\alpha(2)]$,
les trois composantes $M_S = 1, 0, -1$ d'un spin total $S = 1$. La
fonction symétrique
$\Psi_+$ doit être associée à l'unique fonction de spin antisymétrique,
$S = 0$.

\begin{definition}[États singulet et triplet]\label{def:b3:many-electron-atoms:singlet-triplet}
Un état de spin total $S = 0$ est un \emph{état singulet}\index{etat singulet@état singulet}~;
un état de spin total $S = 1$, dont les trois composantes $M_S = 1, 0, -1$
ont la même énergie en l'absence de champ magnétique, est un \emph{état
triplet}\index{etat triplet@état triplet}.
\end{definition}

\begin{definition}[Intégrale d'échange]\label{def:b3:many-electron-atoms:exchange}
Pour deux orbitales $a$ et $b$ et la répulsion électronique $e^2/4\pi\varepsilon_0
r_{12}$, l'\emph{intégrale d'échange}\index{integrale d'echange@intégrale d'échange} est
\[
K_{ab} = \iint a(1)b(2)\,\frac{e^2}{4\pi\varepsilon_0r_{12}}\,b(1)a(2)\,\dd\tau_1\dd\tau_2,
\]
la même intégrale que la répulsion électronique classique $J_{ab}$ (avec
$a(1)b(2)$ des deux côtés), à ceci près que les électrons sont permutés
d'un côté. $K_{ab}$ est positive et n'a pas d'équivalent classique.
\end{definition}

\begin{proposition}[Énergies du singulet et du triplet]\label{prop:b3:many-electron-atoms:helium}
Au premier ordre en la répulsion électronique, les états $1s2s$ de
l'hélium ont pour énergies $E_\pm = E_{1s} + E_{2s} + J \pm K$, le
singulet ($+$) se trouvant au-dessus du triplet ($-$) de $2K$.
\end{proposition}

\begin{proof}
L'énergie non perturbée des deux $\Psi_\pm$ est $E_{1s} + E_{2s}$ (orbitales
hydrogénoïdes de charge nucléaire 2). La correction du premier ordre est
la \omterm{def:b3:quantum-model-systems:hermitian}{valeur moyenne} de la répulsion $\hat V_{12}$ dans l'état normé~:
$\frac12\langle 1s(1)2s(2) \pm 2s(1)1s(2)|\hat V_{12}|1s(1)2s(2) \pm
2s(1)1s(2)\rangle$. Les deux termes directs donnent $\frac12(J + J)$~; les
deux termes croisés donnent $\pm\frac12(K + K)$, puisque $\hat V_{12}$ est
symétrique par rapport aux électrons. Donc $\Delta E = J \pm K$, et les
fonctions de spin, normées et non affectées par $\hat V_{12}$, ne font
que multiplier par 1.
\end{proof}

Le triplet est plus bas bien que le hamiltonien ne contienne pas le
spin~: sa fonction spatiale s'annule quand $\mathbf r_1 = \mathbf r_2$,
de sorte que les deux électrons s'évitent et se repoussent moins. Ce
«~trou de Fermi~» est le contenu physique de la première règle de Hund.

\begin{example}[L'intégrale d'échange de l'hélium, mesurée]\label{ex:b3:many-electron-atoms:K}
% ledger: lev:He.2s3S1, lev:He.2s1S0
Les niveaux $1s2s$ de l'hélium se situent à \qty{159856}{cm^{-1}} (${}^3$S) et
\qty{166277}{cm^{-1}} (${}^1$S) au-dessus de l'état fondamental~: le
singulet est plus haut de
\qty{6421}{cm^{-1}}, soit \qty{0.796}{eV}, et $K(1s,2s) = \qty{0.398}{eV}$.
La figure ci-dessous montre les deux familles.
\end{example}

\section{Couplage de Russell--Saunders et symboles de termes}

Dans un atome léger, les moments cinétiques orbitaux des électrons
s'additionnent en un total $\mathbf L$, leurs spins en un total
$\mathbf S$, et c'est seulement ensuite que
$\mathbf L$ et $\mathbf S$ se couplent faiblement l'un à l'autre.

\begin{definition}[Couplage de Russell--Saunders]\label{def:b3:many-electron-atoms:russell-saunders}
Dans le \emph{couplage de Russell--Saunders}\index{couplage de Russell--Saunders}, les
états d'une configuration sont repérés par le \emph{nombre quantique de
moment cinétique orbital total}\index{nombre quantique de moment cinetique orbital total@nombre quantique de moment cinétique orbital total}
$L$ ($\hat{\mathbf L}^2 = L(L+1)\hbar^2$, avec $M_L = \sum m_l$), le
\emph{nombre quantique de spin total}\index{nombre quantique de spin total} $S$ ($M_S = \sum m_s$),
et le \emph{nombre quantique de moment cinétique total}\index{nombre quantique de moment cinetique total@nombre quantique de moment cinétique total}
$J$, qui prend les valeurs $|L - S|, \dots, L + S$.
\end{definition}

\begin{definition}[Terme spectroscopique, multiplicité de spin, symbole de terme]\label{def:b3:many-electron-atoms:term}
Un \emph{terme spectroscopique}\index{terme spectroscopique} est l'ensemble
des
$(2L+1)(2S+1)$ états d'une configuration de $L$ et $S$ donnés. Sa
\emph{multiplicité de spin}\index{multiplicite de spin@multiplicité de spin} est $2S + 1$. Le \emph{symbole de terme}
\index{symbole de terme} est ${}^{2S+1}L$, avec les lettres S, P, D, F, G, H
pour $L = 0, 1, 2, 3, 4, 5$~; un état de $J$ donné s'écrit ${}^{2S+1}L_J$ ---
par exemple \termsym{3}{P}{2}.
\end{definition}

\begin{proposition}[Couches complètes]\label{prop:b3:many-electron-atoms:closed-shell}
Une sous-couche pleine a $L = 0$ et $S = 0$~; elle ne contribue en rien au
\omterm{def:b3:many-electron-atoms:term}{symbole de terme}.
\end{proposition}

\begin{proof}
Dans une sous-couche pleine, chaque $m_l$ apparaît deux fois et chaque
$m_s = \pm\frac12$ apparaît $2l+1$ fois, donc $M_L = 0$ et $M_S = 0$~; et il
n'y a qu'une façon de la remplir (un seul déterminant). Un état unique
avec $M_L = M_S = 0$ ne peut appartenir qu'à $L = 0$, $S = 0$.
\end{proof}

\begin{proposition}[Dénombrer les états d'une configuration]\label{prop:b3:many-electron-atoms:count}
$N$ électrons dans une sous-couche de $2(2l+1)$ \omterm{def:b3:many-electron-atoms:spin-orbital}{spin-orbitales} donnent
$\binom{2(2l+1)}{N}$ déterminants, et les \omterm{def:b3:quantum-model-systems:degenerate}{dégénérescences} $(2L+1)(2S+1)$
de ses termes ont pour somme ce nombre.
\end{proposition}

\begin{proof}
Un déterminant est fixé par l'ensemble des \omterm{def:b3:many-electron-atoms:spin-orbital}{spin-orbitales} occupées (leur
ordre ne change que son signe)~: on en choisit $N$ parmi les $2(2l+1)$. Les
termes sont les mêmes états regroupés, donc les dénombrements doivent
coïncider.
\end{proof}

\begin{method}[Termes d'une configuration]\label{met:b3:many-electron-atoms:terms}
\begin{enumerate}
  \item Dresser la liste des déterminants permis par le principe de Pauli
        et les ranger dans un tableau selon $M_L$ et $M_S$.
  \item Repérer le plus grand $M_L$~; avec le plus grand $M_S$ qui
        l'accompagne, il ouvre un terme de $L = M_L$, $S = M_S$.
  \item Barrer un déterminant pour chaque couple $(M_L, M_S)$ tel que
        $|M_L| \le L$, $|M_S| \le S$.
  \item Recommencer avec ce qui reste, jusqu'à vider le tableau~;
        vérifier le dénombrement.
\end{enumerate}
\end{method}

\begin{example}[Les termes de $p^2$]\label{ex:b3:many-electron-atoms:p2}
Deux électrons dans $p$ ($m_l = 1, 0, -1$) donnent $\binom62 = 15$
déterminants. Leur tableau selon $M_L$ et $M_S$ est~:
\begin{center}\small
\begin{tabular}{c|ccc}
\hline
 & $M_S = 1$ & $M_S = 0$ & $M_S = -1$\\ \hline
$M_L = 2$ & -- & $(1^+,1^-)$ & --\\
$M_L = 1$ & $(1^+,0^+)$ & $(1^+,0^-)$, $(1^-,0^+)$ & $(1^-,0^-)$\\
$M_L = 0$ & $(1^+,-1^+)$ & $(1^+,-1^-)$, $(1^-,-1^+)$, $(0^+,0^-)$ & $(1^-,-1^-)$\\
$M_L = -1$ & $(0^+,-1^+)$ & $(0^+,-1^-)$, $(0^-,-1^+)$ & $(0^-,-1^-)$\\
$M_L = -2$ & -- & $(-1^+,-1^-)$ & --\\ \hline
\end{tabular}
\end{center}
(avec $m_l^\pm$ pour $m_s = \pm\frac12$). $M_L = 2$ n'apparaît qu'avec $M_S = 0$~:
un terme ${}^1$D (5 états). Le plus grand $M_L = 1$ restant
s'accompagne de
$M_S = 1$~: un terme ${}^3$P (9 états). Il reste un état avec
$M_L = M_S = 0$~: un terme
${}^1$S. On a bien $5 + 9 + 1 = 15$.
\end{example}

\begin{example}[Le carbone, mesuré]\label{ex:b3:many-electron-atoms:carbon}
% ledger: lev:C.3P1, lev:C.3P2, lev:C.1D2, lev:C.1S0
La configuration fondamentale $2p^2$ du carbone donne les niveaux \termsym{3}{P}{0},
\termsym{3}{P}{1}, \termsym{3}{P}{2} à 0, 16.4 et \qty{43.4}{cm^{-1}},
puis
\termsym{1}{D}{2} à \qty{10193}{cm^{-1}} et \termsym{1}{S}{0} à
\qty{21648}{cm^{-1}}~: le terme ${}^3$P est le plus bas, comme le
prévoient les règles de Hund.
\end{example}

\begin{omfigure}
\begin{tikzpicture}
\begin{axis}[name=main, width=0.62\linewidth, height=6.4cm, xmin=-0.3, xmax=6.6,
  ymin=-1500, ymax=23500, axis y line=left, axis x line=none,
  ylabel={énergie / cm$^{-1}$}, unbounded coords=jump, clip=false,
  scaled y ticks=false, ytick={0,10000,20000},
  tick label style={font=\small}, label style={font=\small}]
\addplot[omstate] table[x=x, y=E] {figdata/out/bachelor-3-atomic-levels-carbon.dat};
\draw[densely dotted, gray] (axis cs:1,4858.5) -- (axis cs:2,29.6);
\draw[densely dotted, gray] (axis cs:1,4858.5) -- (axis cs:2,10192.7);
\draw[densely dotted, gray] (axis cs:1,4858.5) -- (axis cs:2,21648);
\draw[densely dotted, gray] (axis cs:3,10192.7) -- (axis cs:4,10192.7);
\draw[densely dotted, gray] (axis cs:3,21648) -- (axis cs:4,21648);
\node[font=\small, anchor=south] at (axis cs:0.5,4858.5) {$2p^2$};
\node[font=\small, anchor=south] at (axis cs:2.5,21648) {${}^1$S};
\node[font=\small, anchor=south] at (axis cs:2.5,10192.7) {${}^1$D};
\node[font=\small, anchor=south] at (axis cs:2.5,29.6) {${}^3$P};
\node[font=\small, anchor=west] at (axis cs:5.05,21648) {\termsym{1}{S}{0}};
\node[font=\small, anchor=west] at (axis cs:5.05,10192.7) {\termsym{1}{D}{2}};
\node[font=\scriptsize, anchor=north] at (axis cs:0.5,-600) {configuration};
\node[font=\scriptsize, anchor=north] at (axis cs:2.5,-600) {termes};
\node[font=\scriptsize, anchor=north] at (axis cs:4.5,-600) {niveaux};
\draw[gray, ->] (axis cs:3.05,60) -- (axis cs:6.9,2500);
\end{axis}
\begin{axis}[at={(main.south east)}, anchor=south west, xshift=1.2cm, yshift=0.9cm,
  width=3.6cm, height=3.6cm, xmin=-0.1, xmax=1.1, ymin=-5, ymax=50,
  axis y line=right, axis x line=none, unbounded coords=jump, ytick={0,16.4,43.4},
  yticklabels={0,16.4,43.4}, tick label style={font=\scriptsize}, clip=false,
  title={\scriptsize ${}^3$P, agrandi}, title style={yshift=-4pt}]
\addplot[omstate] table[x=x, y=E] {figdata/out/bachelor-3-atomic-levels-carbon-zoom.dat};
\node[font=\scriptsize, anchor=east] at (axis cs:0,0) {$J = 0$};
\node[font=\scriptsize, anchor=east] at (axis cs:0,16.4) {$J = 1$};
\node[font=\scriptsize, anchor=east] at (axis cs:0,43.4) {$J = 2$};
\end{axis}
\end{tikzpicture}

{\small La configuration fondamentale du carbone~: une configuration,
trois termes (à leurs énergies mesurées~; le terme ${}^3$P à la moyenne
pondérée de ses niveaux, la configuration à la moyenne des quinze états)
et cinq niveaux. L'éclatement spin--orbite de ${}^3$P (agrandi, en
\unit{cm^{-1}}) est quelques centaines de fois plus petit que les écarts
entre termes.}
\end{omfigure}

\section{Règles de Hund et couplage spin--orbite}

\begin{proposition}[Règles de Hund pour les termes]\label{prop:b3:many-electron-atoms:hund-terms}
Pour la configuration fondamentale d'un atome ou d'un ion, le terme le
plus bas est celui qui a (1) le plus grand $S$, puis (2), pour ce $S$, le
plus grand $L$. Son niveau le plus bas a (3) $J = |L - S|$ si la
sous-couche est moins qu'à moitié pleine, $J = L +
S$ si elle est plus qu'à moitié pleine.
\end{proposition}
\begin{proof}[Statut]
Ces règles sont établies expérimentalement, pour les seules
configurations fondamentales~; ce ne sont pas des théorèmes. La règle 1
est la stabilisation d'échange de la section précédente~; la règle 2 dit
que des électrons qui tournent dans le même sens se rencontrent moins
souvent~; la règle 3 découle du signe du \omterm{def:b3:many-electron-atoms:spin-orbit}{couplage spin--orbite},
ci-dessous.
\end{proof}

\begin{method}[Le terme fondamental à partir des cases quantiques]\label{met:b3:many-electron-atoms:ground-term}
\begin{enumerate}
  \item Remplir les cases de la sous-couche ouverte à partir de
        $m_l = +l$ en descendant, un électron par case avec des spins
        parallèles, puis apparier en repartant de $+l$.
  \item $S = \frac12 \times$(nombre d'électrons célibataires)~; $L = |\sum m_l|$.
  \item $J = |L - S|$ pour une sous-couche moins qu'à moitié pleine,
        $L + S$ pour une sous-couche plus qu'à moitié
        pleine~; pour une sous-couche exactement à moitié pleine,
        $L = 0$ et $J = S$.
\end{enumerate}
\end{method}

\begin{example}[Termes fondamentaux]\label{ex:b3:many-electron-atoms:ground-terms}
% ledger: lev:Ti2.3P0, lev:Ni2.3P0, lev:Cr3.4P1_2
\ce{Ti^2+}, $d^2$~: cases $m_l = 2, 1$ remplies, $S = 1$, $L = 3$, $J = 2$~:
\termsym{3}{F}{2}. \ce{Cr^3+}, $d^3$~: $S = \frac32$, $L = 2 + 1 + 0 = 3$,
\termsym{4}{F}{3/2}. \ce{Fe^2+}, $d^6$~: cinq spins parallèles, plus un spin
opposé dans $m_l = 2$~:
$S = 2$, $L = 2$, plus qu'à moitié pleine, \termsym{5}{D}{4}. \ce{Ni^2+}, $d^8$~:
$S = 1$, $L = 3$, \termsym{3}{F}{4}. Chacun s'accorde avec le niveau
fondamental que mesure la spectroscopie atomique.
\end{example}

\begin{definition}[Couplage spin--orbite, structure fine]\label{def:b3:many-electron-atoms:spin-orbit}
Le \emph{couplage spin--orbite}\index{couplage spin--orbite} est
l'interaction entre les moments magnétiques associés au spin d'un
électron et à son mouvement orbital, écrite $A\,\hat{\mathbf L}\cdot\hat{\mathbf S}/\hbar^2$
pour un terme. Il éclate un terme en ses niveaux de $J$ différents~: c'est
la \emph{structure fine}\index{structure fine} du terme.
\end{definition}

\begin{proposition}[Règle des intervalles de Landé]\label{prop:b3:many-electron-atoms:lande}
Au sein d'un terme, $E(J) = E_0 + \frac{A}{2}[J(J+1) - L(L+1) - S(S+1)]$, de sorte que
$E(J) - E(J - 1) = AJ$. $A > 0$ (structure normale) pour une sous-couche
moins qu'à moitié pleine,
$A < 0$ (structure inversée) pour une sous-couche plus qu'à moitié pleine.
\end{proposition}

\begin{proof}
$\hat{\mathbf J} = \hat{\mathbf L} + \hat{\mathbf S}$ donne $\hat{\mathbf J}^2 =
\hat{\mathbf L}^2 + \hat{\mathbf S}^2 + 2\hat{\mathbf L}\cdot\hat{\mathbf S}$, donc
$\hat{\mathbf L}\cdot\hat{\mathbf S} = \frac12(\hat{\mathbf J}^2 - \hat{\mathbf L}^2 -
\hat{\mathbf S}^2)$, dont la \omterm{def:b3:quantum-model-systems:operator}{valeur propre} dans un état de $J$, $L$, $S$
donnés est
$\frac{\hbar^2}{2}[J(J+1) - L(L+1) - S(S+1)]$. La différence entre $J$ et
$J - 1$ vaut $\frac A2[J(J+1) - (J-1)J] = AJ$. Le signe de $A$ est admis~:
une sous-couche plus qu'à moitié pleine se comporte comme le nombre
correspondant de trous positifs.
\end{proof}

\begin{example}[Mise à l'épreuve de la règle des intervalles]\label{ex:b3:many-electron-atoms:lande-test}
% ledger: lev:C.3P1, lev:C.3P2, lev:O.3P1, lev:O.3P0
Pour ${}^3$P, la règle prévoit des intervalles dans le rapport $J = 2 : 1$,
soit 2. Carbone~: $(43.4 - 16.4)/16.4 = 1.65$. Oxygène ($2p^4$, structure
inversée, \termsym{3}{P}{2}
le plus bas, puis $J = 1$ à \qty{158.3}{cm^{-1}} et $J = 0$ à \qty{227.0}{cm^{-1}})~:
$158.3/68.7 = 2.30$. La règle tient à 20\,\% près~: le \omterm{def:b3:many-electron-atoms:russell-saunders}{couplage de
Russell--Saunders} décrit bien les atomes légers, sans être exact.
\end{example}

\begin{omfigure}
\begin{tikzpicture}[font=\small]
  \draw[omstate] (0,0) -- (1.6,0) node[midway, below] {$3s\ {}^2$S$_{1/2}$};
  \draw[omstate] (0,3.0) -- (1.6,3.0);
  \node[anchor=east] at (-0.1,3.0) {$3p$};
  \draw[omstate] (2.8,2.75) -- (4.4,2.75) node[right] {${}^2$P$_{1/2}$};
  \draw[omstate] (2.8,3.25) -- (4.4,3.25) node[right] {${}^2$P$_{3/2}$};
  \draw[densely dotted, gray] (1.6,3.0) -- (2.8,2.75) (1.6,3.0) -- (2.8,3.25);
  \draw[omfluo] (3.3,2.73) -- (3.3,0.02) node[pos=0.55, left] {D$_1$};
  \draw[omfluo] (3.9,3.23) -- (3.9,0.02) node[pos=0.55, right] {D$_2$};
  \draw[omstate] (2.8,0) -- (4.4,0);
  \draw[densely dotted, gray] (1.6,0) -- (2.8,0);
  \draw[decorate, decoration={brace, amplitude=3pt}] (5.5,3.27) -- (5.5,2.73)
     node[midway, right=3pt, align=left] {\footnotesize \qty{17.2}{cm^{-1}}\\[-1pt]\footnotesize (pas à l'échelle)};
\end{tikzpicture}

{\small Les raies D du sodium. Le niveau $3p$ est éclaté par le \omterm{def:b3:many-electron-atoms:spin-orbit}{couplage
spin--orbite} en ${}^2$P$_{1/2}$ et ${}^2$P$_{3/2}$~; l'émission vers le
niveau fondamental $3s$ donne deux raies jaunes, D$_1$ à
\qty{589.76}{nm} et D$_2$ à \qty{589.16}{nm} (longueurs d'onde dans le
vide).}
\end{omfigure}

\begin{example}[Le doublet du sodium]\label{ex:b3:many-electron-atoms:sodium}
% ledger: lev:Na.3p1_2, lev:Na.3p3_2
Les niveaux $3p$ du sodium se situent à 16956.17 et
\qty{16973.37}{cm^{-1}}~: les raies D sont à
$10^7/16956.17 = \qty{589.76}{nm}$ et \qty{589.16}{nm} dans le vide,
séparées de \qty{17.20}{cm^{-1}}. Pour ${}^2$P, $E(\frac32) - E(\frac12) =
\frac32A$, donc $A = \qty{11.47}{cm^{-1}}$.
\end{example}

\begin{proposition}[Règles de sélection des atomes]\label{prop:b3:many-electron-atoms:selection-atoms}
Dans le \omterm{def:b3:many-electron-atoms:russell-saunders}{couplage de Russell--Saunders}, une transition dipolaire
électrique exige
$\Delta S = 0$, $\Delta L = 0, \pm1$, $\Delta J = 0, \pm1$ (mais pas $J = 0 \to J
= 0$), et un changement de parité~: pour le saut d'un seul électron,
$\Delta l = \pm1$.
\end{proposition}
\begin{proof}[Statut]
$\Delta S = 0$ est démontré au \cref{ch:b3:electronic-spectroscopy}~:
l'\omterm{def:b3:quantum-model-systems:operator}{opérateur} dipolaire n'agit pas sur le spin. Les autres règles découlent
du caractère vectoriel de l'\omterm{def:b3:quantum-model-systems:operator}{opérateur} dipolaire et sont admises ici.
\end{proof}

\begin{omfigure}
\begin{tikzpicture}
\begin{axis}[width=0.62\linewidth, height=6.2cm, xmin=-0.5, xmax=3.5,
  ymin=158, ymax=173, axis y line=left, axis x line=none,
  ylabel={énergie / $10^3$ cm$^{-1}$}, unbounded coords=jump, clip=false,
  tick label style={font=\small}, label style={font=\small}]
\addplot[omstate] table[x=x, y=E] {figdata/out/bachelor-3-atomic-levels-helium.dat};
\node[font=\small, anchor=west] at (axis cs:1.05,166.277) {$1s2s\ {}^1$S};
\node[font=\small, anchor=west] at (axis cs:1.05,171.135) {$1s2p\ {}^1$P};
\node[font=\small, anchor=west] at (axis cs:3.05,159.856) {$1s2s\ {}^3$S};
\node[font=\small, anchor=west] at (axis cs:3.05,169.087) {$1s2p\ {}^3$P};
\node[font=\small, anchor=south] at (axis cs:0.5,172.3) {singulets};
\node[font=\small, anchor=south] at (axis cs:2.5,172.3) {triplets};
\draw[omfluo] (axis cs:0.5,171.0) -- (axis cs:0.5,166.42) node[pos=0.5, left, font=\scriptsize] {\qty{2059}{nm}};
\draw[omfluo] (axis cs:2.5,168.95) -- (axis cs:2.5,160.0) node[pos=0.5, left, font=\scriptsize] {\qty{1083}{nm}};
\draw[densely dotted, gray] (axis cs:-0.3,159.856) -- (axis cs:2,159.856);
\draw[<->, omThm] (axis cs:-0.2,159.856) -- (axis cs:-0.2,166.277) node[pos=0.5, right, font=\scriptsize] {$2K$};
\end{axis}
\end{tikzpicture}

{\small Les deux «~héliums~». Niveaux singulets et triplets des
configurations $1s2s$ et $1s2p$ à leurs énergies mesurées au-dessus de
l'état fondamental (le terme ${}^3$P à la moyenne de ses niveaux). Les
traits ne relient que des états de la même famille~; chaque triplet se
trouve sous son singulet, à $2K$ au-dessous.}
\end{omfigure}

\begin{history}[Deux héliums et le principe d'exclusion]
L'hélium fut nommé d'après le Soleil en 1868 et découvert sur Terre par
William Ramsay en 1895. Ses deux séries de raies déroutèrent les
spectroscopistes pendant trente ans. En 1925, Wolfgang Pauli proposa
que deux électrons ne puissent pas partager les mêmes nombres
quantiques, l'année même où George Uhlenbeck et Samuel Goudsmit
proposaient le spin de l'électron~; en 1926, Werner Heisenberg montra que
la seule symétrie de la fonction d'onde scinde l'hélium en ses familles
singulet et triplet.
\end{history}

\begin{inthelab}[Résoudre le doublet du sodium]
Une lampe à sodium est placée devant la fente d'un spectromètre à réseau.
Avec un réseau de 600 traits par millimètre, les angles du premier ordre
des deux raies diffèrent d'environ \num{0.02}\textdegree~: un bon
spectromètre les sépare, un prisme d'étudiant non. La lampe est très
chaude~; on ne la manipule qu'une fois refroidie.
\end{inthelab}

\section{Exercices}

\begin{exercise}[$\star$]\label{exo:b3:many-electron-atoms:1}
Écrire le \omterm{def:b3:many-electron-atoms:slater}{déterminant de Slater} de l'état fondamental du lithium,
$1s^22s$, l'électron $2s$ ayant le spin $\alpha$. Pourquoi n'existe-t-il
pas de déterminant pour
$1s^3$~?
\end{exercise}

\begin{exercise}[$\star$]\label{exo:b3:many-electron-atoms:2}
Dénombrer les déterminants des configurations $p^2$, $p^3$ et $d^2$. Les
termes de $p^3$ sont ${}^4$S, ${}^2$D, ${}^2$P et ceux de $d^2$ sont ${}^3$F, ${}^3$P, ${}^1$G,
${}^1$D, ${}^1$S~: vérifier les deux dénombrements.
\end{exercise}

\begin{exercise}[$\star$]\label{exo:b3:many-electron-atoms:3}
Donner le terme et le niveau fondamentaux de N, O, F, \ce{Ti^2+} et \ce{Fe^3+} ($d^5$).
\end{exercise}

\begin{exercise}[$\star$]\label{exo:b3:many-electron-atoms:4}
Énumérer les niveaux des termes ${}^2$D et ${}^3$P avec leurs
\omterm{def:b3:quantum-model-systems:degenerate}{dégénérescences} $2J+1$, et vérifier que la somme des \omterm{def:b3:quantum-model-systems:degenerate}{dégénérescences} vaut
$(2L+1)(2S+1)$.
\end{exercise}

\begin{exercise}[$\star\star$]\label{exo:b3:many-electron-atoms:5}
Les niveaux ${}^3$P du carbone se situent à 0, 16.4 et
\qty{43.4}{cm^{-1}}. Calculer le rapport des intervalles et la constante
de \omterm{def:b3:many-electron-atoms:spin-orbit}{couplage spin--orbite} $A$ tirée de chaque intervalle. Que traduit
l'écart entre les deux valeurs de $A$~?
\end{exercise}

\begin{exercise}[$\star\star$]\label{exo:b3:many-electron-atoms:6}
Les niveaux $3p$ du sodium se situent à 16956.17 et
\qty{16973.37}{cm^{-1}}. Calculer les longueurs d'onde dans le vide des
raies D, leur écart en \unit{cm^{-1}} et en meV, et $A$ pour le terme $3p$.
\end{exercise}

\begin{exercise}[$\star\star$]\label{exo:b3:many-electron-atoms:7}
Parmi ces transitions du sodium, lesquelles sont permises en émission~:
$3p \to 3s$,
$3d \to 3s$, $3d \to 3p$, $4s \to 3s$, $4s \to 3p$~? Justifier chaque
réponse.
\end{exercise}

\begin{exercise}[$\star\star$]\label{exo:b3:many-electron-atoms:8}
% ledger: lev:He.2p3P2, lev:He.2p3P1, lev:He.2p3P0, lev:He.2p1P1
Les niveaux $1s2p$ de l'hélium sont \termsym{3}{P}{2}, \termsym{3}{P}{1},
\termsym{3}{P}{0} à 169086.76, 169086.84, \qty{169087.83}{cm^{-1}} et
\termsym{1}{P}{1} à \qty{171134.89}{cm^{-1}}. Calculer l'énergie moyenne du
terme ${}^3$P, puis $K(1s,2p)$ en eV. Comparer à $K(1s,2s) = \qty{0.398}{eV}$
et expliquer la différence.
\end{exercise}

\begin{exercise}[$\star\star$]\label{exo:b3:many-electron-atoms:9}
Trouver les termes de $d^2$ par la méthode de ce chapitre, en partant du
plus grand $M_L$ (on peut se dispenser d'écrire tout le tableau, mais il
faut justifier chaque terme).
\end{exercise}

\begin{exercise}[$\star\star\star$]\label{exo:b3:many-electron-atoms:10}
Montrer que la fonction spatiale triplet $\frac{1}{\sqrt2}[1s(1)2s(2) -
2s(1)1s(2)]$ s'annule quand $\mathbf r_1 = \mathbf r_2$ et que la fonction
singulet ne s'annule pas. Relier ce résultat au signe de $E_+ - E_-$.
\end{exercise}

\begin{exercise}[$\star\star\star$]\label{exo:b3:many-electron-atoms:11}
% ledger: lev:Ni2.3F3, lev:Ni2.3F2
Pour \ce{Ni^2+} ($d^8$), les niveaux ${}^3$F se situent à 0 ($J = 4$),
1360.7 ($J = 3$)
et \qty{2269.6}{cm^{-1}} ($J = 2$). Expliquer cet ordre, calculer le
rapport des intervalles et le comparer à la prévision de Landé. Donner
$A$ à partir du plus grand intervalle.
\end{exercise}

\begin{exercise}[$\star\star\star$]\label{exo:b3:many-electron-atoms:12}
Montrer que la \omterm{def:b3:quantum-model-systems:hermitian}{valeur moyenne} de $\hat{\mathbf L}\cdot\hat{\mathbf S}$,
sommée sur l'ensemble des $(2L+1)(2S+1)$ états d'un terme, est nulle, de
sorte que le \omterm{def:b3:many-electron-atoms:spin-orbit}{couplage spin--orbite} laisse inchangée la moyenne pondérée
des niveaux. Le vérifier sur le terme ${}^3$P du carbone.
\end{exercise}

\section{Problème~: deux héliums}

\begin{problem}[{Problème du week-end --- les familles singulet et
triplet de l'hélium~: déterminants, termes des ions des métaux de
transition, règles de sélection qui séparent les deux familles, et
intégrale d'échange mesurée}]
\label{pb:b3:many-electron-atoms:1}
% ledger: lev:He.2s3S1, lev:He.2s1S0, lev:He.2p3P2, lev:He.2p3P1, lev:He.2p3P0, lev:He.2p1P1, lev:Ti2.3F3, lev:Ti2.3F4, lev:Ti2.3P0, lev:Ti2.3P1, lev:Ti2.3P2
Données (NIST, en \unit{cm^{-1}} au-dessus de l'état fondamental de
chaque espèce). Hélium~:
$1s2s$ \termsym{3}{S}{1} 159855.97, \termsym{1}{S}{0} 166277.44~; $1s2p$
\termsym{3}{P}{2,1,0} 169086.76, 169086.84, 169087.83, \termsym{1}{P}{1}
171134.89. \ce{Ti^2+}~: \termsym{3}{F}{2,3,4} 0, 184.9, 420.4~; \termsym{3}{P}{0,1,2}
10538.4, 10603.6, 10721.2. $\qty{1}{eV} = \qty{8065.54}{cm^{-1}}$.

\medskip\noindent\textbf{Partie I --- \omterm{def:b3:many-electron-atoms:spin-orbital}{Spin-orbitales} et déterminants.}
\begin{enumerate}
  \item Énumérer les quatre \omterm{def:b3:many-electron-atoms:spin-orbital}{spin-orbitales} disponibles pour la
        configuration $1s2s$.
  \item Écrire les quatre déterminants ayant un électron dans $1s$ et un
        dans $2s$.
  \item Montrer que $|1s\alpha\,2s\alpha|$ est le produit d'une fonction
        spatiale antisymétrique et de la fonction de spin $\alpha(1)\alpha(2)$.
  \item Combiner $|1s\alpha\,2s\beta|$ et $|1s\beta\,2s\alpha|$ pour obtenir
        la composante $M_S = 0$ de la même fonction spatiale, ainsi qu'un
        quatrième état.
  \item Quelle fonction spatiale, symétrique ou antisymétrique, va avec
        $S = 0$~? Avec $S = 1$~?
  \item Pourquoi parle-t-on de singulet et de triplet~?
\end{enumerate}

\medskip\noindent\textbf{Partie II --- Les termes de \ce{Ti^2+}.}
\begin{enumerate}[resume]
  \item Donner la configuration de \ce{Ti^2+} et le nombre de ses
        déterminants.
  \item Trouver son terme et son niveau fondamentaux à l'aide des règles
        de Hund~; les confronter aux données.
  \item Ses autres termes sont ${}^3$P, ${}^1$G, ${}^1$D et ${}^1$S~:
        vérifier le dénombrement des états.
  \item La \omterm{def:b3:many-electron-atoms:spin-orbit}{structure fine} de ${}^3$F est-elle normale ou inversée~?
        Mettre à l'épreuve la règle des intervalles de Landé.
  \item Calculer les énergies moyennes pondérées des termes ${}^3$F et
        ${}^3$P.
  \item Leur écart vaut $15B$, où $B$ est un paramètre de Racah de l'ion
        (\cref{ch:b3:complex-spectra-magnetism}). Calculer $B$.
\end{enumerate}

\medskip\noindent\textbf{Partie III --- Pourquoi deux familles.}
\begin{enumerate}[resume]
  \item Quelle règle de sélection interdit une transition entre un
        singulet et un triplet~?
  \item Le triplet le plus bas, $1s2s$ \termsym{3}{S}{1}, ne peut pas
        émettre vers l'état fondamental $1s^2$ \termsym{1}{S}{0}. En donner
        deux raisons. Comment appelle-t-on un tel état~?
  \item Calculer la longueur d'onde de la raie $1s2p\ {}^3$P $\to 1s2s\ {}^3$S (utiliser
        la moyenne des niveaux ${}^3$P).
  \item Calculer la longueur d'onde de $1s2p\ {}^1$P $\to 1s2s\ {}^1$S.
  \item Calculer la longueur d'onde de $1s2p\ {}^1$P $\to 1s^2\ {}^1$S.
        Dans quel domaine du spectre se trouve-t-elle~?
  \item Expliquer pourquoi les spectroscopistes du \textsc{xix}\textsuperscript{e}
        siècle, voyant ces raies, crurent à deux gaz différents.
\end{enumerate}

\medskip\noindent\textbf{Partie IV --- L'\omterm{def:b3:many-electron-atoms:exchange}{intégrale d'échange}, mesurée.}
\begin{enumerate}[resume]
  \item Écrire les énergies au premier ordre du singulet et du triplet
        $1s2s$ en fonction de $E_{1s}$, $E_{2s}$, $J$ et $K$.
  \item Lequel est le plus bas, et pourquoi, en termes de positions des
        électrons~?
  \item Calculer l'écart singulet--triplet de $1s2s$ en \unit{cm^{-1}} et
        en eV.
  \item Calculer $K(1s,2p)$ en eV à partir des niveaux $1s2p$.
  \item Pourquoi $K(1s,2s)$ est-elle plus grande que $K(1s,2p)$~?
  \item La première règle de Hund est-elle confirmée par les états
        excités de l'hélium~?
  \item Énoncer le résultat~: l'\omterm{def:b3:many-electron-atoms:exchange}{intégrale d'échange} $K(1s,2s)$ de
        l'hélium, en eV.
\end{enumerate}
\end{problem}
